\subsection{Exact family transport with necessary order filters}
Let $A\leq\operatorname{Aut}(G)$ be a finite, explicitly computed
transformation group, and let $\mathcal H=(H_1,\ldots,H_n)$ and
$\mathcal K=(K_1,\ldots,K_n)$ be ordered families. We seek an element of
$A$ carrying the first family to a permutation of the second. Put
$H_J=\bigcap_{j\in J}H_j$, with $H_\varnothing=G$, and define $K_J$
similarly. If a proposed type permutation sends $H_i$ to $K_{\pi(i)}$,
then necessarily
\[
 |H_J|=|K_{\pi(J)}|\qquad(J\subseteq\{1,\ldots,n\}).       \tag{**}
\]
Indeed an automorphism commutes with intersection and preserves order.
Thus a failed equality rejects that type permutation, and a failed
comparison of the two multisets of intersection orders rejects every
permutation. Satisfying (**) never certifies equivalence.

Fix a surviving target tuple $\mathcal T$ and any order
$i_1,\ldots,i_n$ in which to place its entries. Start with $C=A$ and
$\mathcal H' =\mathcal H$. At step $k$ find an actual transporter
$x\in C$ with $x(H'_{i_k})=T_{i_k}$. If none exists, reject this target.
Otherwise replace every $H'_i$ by $x(H'_i)$ and replace $C$ by the
setwise stabilizer of $T_{i_k}$ in $C$. After the last step no further
stabilizer is needed. Success supplies a product of the computed
automorphisms taking the entire family to the target.

For completeness, after the first $k-1$ placements all possible further
solutions differ from the current partial transporter by an element of
the pointwise stabilizer of the already placed \emph{subgroups}. If $x$
is one solution to the next placement, all such solutions differ from
$x$ by an element fixing its target subgroup. Hence passing to the stated
stabilizer loses no solution. Induction proves both soundness and
completeness, for every placement order. The stabilizer after $k$
placements depends only on $A$ and the first $k$ target subgroups, not on
the input family or the transporters chosen. It can therefore be cached
for repeated comparisons against the same target and placement order.
In this implementation larger target subgroups are placed first, with
indices breaking ties. This changes computation time, not the criterion.
Literal equality of complete tuples is an immediate successful case.

\subsection{Reducing automorphisms to simultaneous conjugacy}
There is a useful version when a complete set $B$ of representatives
for $\operatorname{Aut}(G)/\operatorname{Inn}(G)$ is available. For each
$\beta\in B$ and each target type permutation, apply $\beta$ to the
whole input family. The remaining question is simultaneous conjugacy
by one element of $G$. In the preceding procedure start with $C=G$;
replace subgroup transport under automorphisms by conjugation, and
replace each target stabilizer by its normalizer in $C$.

Every automorphism is an inner automorphism composed with one tested
representative, since the inner automorphisms form a normal subgroup.
For a fixed representative, the same transporter-coset argument proves
that the successive normalizers preserve exactly all possible remaining
conjugators. A positive result records the representative's images of
the ambient generators, one actual conjugating element, and the type
permutation; together these identify the entire family. Neither a
splitting of the automorphism group nor a representative of the same
order as its quotient image is required. In particular, when the inner
subgroup has index two, any verified noninner automorphism together
with the identity is a sufficient complete transversal. This is a
standard quotient and transporter reduction; its implementation below
is not a new classification theorem.

